\subsection{反模}

定义 3. --- 设 A 为环，M 为 A-模，\(E = \text{End}_{A}(M)\) 为 M 的自同态环。M 的反模是与 M 具有相同底层加群、外律为 \((c, x) \mapsto c(x)\) 的左 E-模。

设 Z 为环 A 的中心。对每个 \(a \in Z\) ，位似 \(a_{M}\) 属于 E。因此，E 典范地带有 Z-代数结构。特别地，若 M 是有限生成的 Z-模，则 M 的反模是有限生成的。

引理 4. --- 设 M 为反模有限生成的左 A-模。则存在自然数 m 以及一个单 \(A_{M}\) -线性映射，它从 \((\mathrm{A}_{\mathrm{M}})_{s}\) 映到 \(M^{m}\) 。

置 \(\mathrm{E}=\mathrm{End}_{\mathrm{A}}(\mathrm{M})\) 。设 \((x_{1},\ldots,x_{m})\) 为 E-模 M 的有限生成族。映射 \(\varphi:a\mapsto(ax_{1},\ldots,ax_{m})\) 把 \((\mathrm{A}_{\mathrm{M}})_{s}\) 映到 \(M^{m}\) ，它是 \(A_{M}\) -线性的。设 a 为 \(A_{M}\) 的元素且 \(\varphi(a)=0\) 。M 中使 ax = 0 的元素 x 组成的集合是 M 的包含 \(x_{1}, \ldots, x_{m}\) 的 E-子模，因此等于 M，由此推出 a = 0。

命题 6. --- 设 M 为反模有限生成的阿廷（相应地，诺特）左 A-模。则 M 的位似环 \(A_{M}\) 是左阿廷（相应地，左诺特）的。

这由引理 4 及 VIII, p.~3 的命题 3 得出。

推论. --- 设 A 为交换环。

\begin{enumerate}
\def\labelenumi{\alph{enumi})}
\item
  设 M 为诺特 A-模。则环 \(A_{M}\) 是诺特的。
\item
  设 M 为有限长度的 A-模。则环 \(A_{M}\) 是阿廷的。
\end{enumerate}

设 M 为 A-模。在 a) 或 b) 的假设下，A-模 M 是有限生成的。由于 A 交换，\(A_{M}\) 含于环 \(\operatorname{End}_{\mathrm{A}}(\mathrm{M})\) ，故 M 的反模是有限生成的。于是只需应用命题 6。

注. --- 设 A 为环。反模有限生成的阿廷左 A-模具有有限长度：事实上，M 的位似环 \(A_{M}\) 是左阿廷的（命题 6），而 M 是 \(A_{M}\) 上的阿廷模；由 VIII, p.~7 的命题 4，模 M 在 \(A_{M}\) 上具有有限长度，因而在 A 上亦然。

特别地，交换环上的每个有限生成阿廷模都具有有限长度。与此相对照，非交换环上的有限生成阿廷模不一定具有有限长度（VIII, p.~16, 习题 12）。

